\subsection{The Language of Numbers and the Power of Names}
\label{sec:numbers-and-names}
\slidesrc{Pragmatics: 50--66}

\subsubsection{A few words on maths}
\slidesrc{Pragmatics: 50--59}

Galileo Galilei (1564--1642), in \emph{Il Saggiatore}, called mathematics a language:

\begin{quote}\small
``Philosophy is written in this grand book, which stands continually open before our
eyes (I say the `Universe'), but cannot be understood without first learning to
comprehend the language and know the characters in which it is written. It is written in
mathematical language, and its characters are triangles, circles and other geometric
figures, without which it is impossible to humanly understand a word; without these one
is wandering in a dark labyrinth.''
\end{quote}

If mathematics is a language, the way numbers are \emph{written and named} should also
matter. The slides follow its history:
\begin{itemize}
  \item \textbf{Greek numerals}. The ancient Greeks first had a system like the Romans;
        from the 4th century BC they used different symbols for 1--9 (as we do), but
        \emph{another} set for 10, 20, 30, \ldots\ and yet another for 100, 200, 300:
        like many ancient systems, they eventually ran out of symbols. They had \textbf{no
        symbol for zero} (\emph{mēdén}) and no digit placeholder: they were unsure
        whether zero was a number at all (``how can nothing be something?''). Yet with
        these numerals \textbf{Eratosthenes} of Cyrene (c.\ 276--194 BC), chief
        librarian of Alexandria, computed the circumference of the Earth with remarkable
        accuracy, invented geography as a discipline, and gave the \emph{sieve of
        Eratosthenes} for finding prime numbers.
  \item \textbf{Roman numerals}. Seven symbols: I $=1$, V $=5$, X $=10$, L $=50$, C
        $=100$, D $=500$, M $=1000$. They remained the usual way of writing numbers in
        Europe well into the Late Middle Ages and were replaced from the 14th century by
        the more convenient Hindu-Arabic numerals; they survive in minor uses today. The
        Romans nevertheless built aqueducts and roads that still stand.
  \item \textbf{Hindu-Arabic numerals}: ten symbols, 1, 2, \ldots, 9, 0, for the
        decimal positional system. They originated in India in the 6th or 7th century
        and reached Europe through Middle Eastern mathematicians (especially
        al-Khwarizmi and al-Kindi) around the 12th century. They were a profound break
        with earlier methods of counting (such as the abacus) and paved the way for
        algebra.
  \item \textbf{Counting with the body}. The Oksapmin of Papua New Guinea count with a
        system of 27 body parts: from the thumb of one hand, up the arm, across the head
        and down the other arm; the number is the name of the body part reached.
  \item \textbf{Languages with few number words}. Some Amazonian languages have no words
        for exact numbers beyond very small ones. In matching tasks (put on the table as
        many objects as you see in a row), speakers succeed with small quantities but
        become approximate with larger ones: without number words, exact large
        quantities are hard to keep in mind.
\end{itemize}

\begin{nbintuition}[the message]
Notation and number words are a language, and some languages make some thoughts easier:
long multiplication is easy in positional decimal notation and painful with Roman
numerals; exact counting is hard without number words. But the Greeks and Romans did
great mathematics and engineering with clumsy numerals. Again: language \emph{influences}
thought (weak Whorf), it does not \emph{determine} it.
\end{nbintuition}

\subsubsection{The importance of what a thing is called}
\slidesrc{Pragmatics: 60--66}

The same thing can be named in ways that push the hearer to judge it differently:
\begin{itemize}[nosep]
  \item \emph{patriots}, \emph{activists}, or \emph{terrorists}?
  \item \emph{torture} or \emph{enhanced interrogation techniques}?
\end{itemize}

\textbf{But sometimes language seems ineffective. Why?} In 2003 the cafeterias of the US
House of Representatives renamed French fries \textbf{``freedom fries''} and French toast
``freedom toast'', a symbolic protest against France's opposition to the war in Iraq
(BBC News, 12 March 2003). It had happened before: during World War~I sauerkraut became
``liberty cabbage'', liverwurst ``liberty sausage'', hamburger ``Salisbury steak'',
dachshunds ``liberty pups''. In 2003 the joke list went on: freedom poodles, freedom
kissing, freedom manicures, and ``France'' itself would become ``freedom land'', ``French''
``the language of freedom''. These renamings did not change much: everybody knew what
the new names stood for, so the old meaning (and the old attitude) remained.

\textbf{And sometimes it works.} Prunes had an image problem (old people, laxatives).
Producers renamed them \textbf{dried plums}, and sales went up: ``Folks favor dried plums
over prunes'', with sales up 5.5\% over the previous year, according to the state
marketing board (Business Journal, 2001). Many other examples from all around the world
concern the names of products, and the classic marketing mistakes when a name has an
unwanted meaning in another language.

\begin{nbwhy}
Naming is pragmatics in action: choosing a word is choosing what the hearer will notice
(an agent or not, an act of violence or a technique, an old person's remedy or a fruit).
NLP systems meet this every day: sentiment and stance detection must recognise that
\emph{terrorist} and \emph{freedom fighter} may refer to the same person with opposite
attitudes, and text generation systems choose words with a pragmatic effect, wanted or
not.
\end{nbwhy}

The deck ends with Nelson Mandela:

\begin{quote}
``If you talk to a man in a language he understands, that goes to his head. If you talk
to him in his language, that goes to his heart.''
\end{quote}

\begin{nbsummary}[pragmatics]
\begin{itemize}[nosep]
  \item Pragmatics: how context (situation, pre-existing knowledge, inferred intent)
        contributes to meaning. Semiotics: the study of signs and meaning-making.
  \item Communication: intentional exchange of information through a shared system of
        signs. Classical view (sentences are true/false) vs.\ modern view (language is
        action).
  \item Wittgenstein (1953): meaning is use in a language-game. Austin (1962): in saying
        something we do something (performatives). Searle (1969): speech acts,
        illocutions, intentionality.
  \item Speech acts: inform, query, command, promise, acknowledge. Stages: intention,
        generation, synthesis (speaker); perception, analysis, disambiguation,
        incorporation (hearer).
  \item Language and thought: Sapir--Whorf strong (false) vs.\ weak (supported):
        time/space, colour, gender, agency, numbers.
  \item Names matter: freedom fries (ineffective), dried plums (effective).
\end{itemize}
\end{nbsummary}
